\documentclass[11pt,a4paper]{article}

% ── Packages ──────────────────────────────────────────────────────────────────
\usepackage[utf8]{inputenc}
\usepackage[T1]{fontenc}
\usepackage[french]{babel}
\usepackage[margin=2cm, top=1.8cm, bottom=2cm]{geometry}
\usepackage{graphicx}
\usepackage{booktabs}
\usepackage{array}
\usepackage{tabularx}
\usepackage{fancyhdr}
\usepackage{mdframed}
\usepackage{listings}
\usepackage{xcolor}
\usepackage{amssymb}
\usepackage{enumitem}
\usepackage{tikz}
\usepackage{lmodern}
\usepackage{microtype}
\usepackage{hyperref}
\usepackage{multicol}

% ── Colors ────────────────────────────────────────────────────────────────────
\definecolor{uirblue}{RGB}{0,83,155}
\definecolor{uirgold}{RGB}{212,160,23}
\definecolor{codebg}{RGB}{245,245,245}
\definecolor{codeframe}{RGB}{200,200,200}
\definecolor{exbg}{RGB}{248,249,252}
\definecolor{exframe}{RGB}{0,83,155}
\definecolor{answerbg}{RGB}{255,255,255}
\definecolor{answerframe}{RGB}{80,80,80}
\definecolor{qcmbg}{RGB}{252,252,252}

% ── Code listings ─────────────────────────────────────────────────────────────
\lstdefinelanguage{PHP2}{
  language=PHP,
  morekeywords={class,extends,implements,abstract,interface,public,private,
                protected,function,return,new,null,true,false,static,parent,
                readonly,void,bool,int,float,string,array},
  keywordstyle=\color{uirblue}\bfseries,
  stringstyle=\color{red!70!black},
  commentstyle=\color{gray}\itshape,
  basicstyle=\ttfamily\footnotesize,
  backgroundcolor=\color{codebg},
  frame=single,
  framerule=0.4pt,
  rulecolor=\color{codeframe},
  showstringspaces=false,
  breaklines=true,
  numbers=left,
  numberstyle=\tiny\color{gray},
  numbersep=6pt,
  xleftmargin=14pt,
  tabsize=4,
}

\lstdefinelanguage{HTML2}{
  language=HTML,
  basicstyle=\ttfamily\footnotesize,
  backgroundcolor=\color{codebg},
  frame=single,
  framerule=0.4pt,
  rulecolor=\color{codeframe},
  showstringspaces=false,
  breaklines=true,
  numbers=left,
  numberstyle=\tiny\color{gray},
  numbersep=6pt,
  xleftmargin=14pt,
  tabsize=4,
  keywordstyle=\color{uirblue}\bfseries,
  stringstyle=\color{red!70!black},
}

% ── Answer box ────────────────────────────────────────────────────────────────
\newcommand{\answerbox}[1]{%
  \begin{mdframed}[
    linecolor=answerframe,
    linewidth=0.5pt,
    backgroundcolor=answerbg,
    innertopmargin=6pt,
    innerbottommargin=#1,
    innerleftmargin=6pt,
    innerrightmargin=6pt,
    skipabove=4pt,
    skipbelow=6pt,
  ]
  \phantom{x}
  \end{mdframed}
}

% ── Exercise frame ────────────────────────────────────────────────────────────
\mdfdefinestyle{exstyle}{
  linecolor=exframe,
  linewidth=1.2pt,
  backgroundcolor=exbg,
  innertopmargin=8pt,
  innerbottommargin=6pt,
  innerleftmargin=10pt,
  innerrightmargin=10pt,
  skipabove=10pt,
  skipbelow=4pt,
  frametitlebackgroundcolor=uirblue,
  frametitlefont=\color{white}\bfseries\normalsize,
}

% ── QCM checkbox ─────────────────────────────────────────────────────────────
\newcommand{\cb}{$\square$\ }

% ── Header / footer ──────────────────────────────────────────────────────────
\pagestyle{fancy}
\fancyhf{}
\renewcommand{\headrulewidth}{0.4pt}
\fancyhead[L]{\small\color{gray} Examen Final — Web Development I}
\fancyhead[R]{\small\color{gray} ESIN Ing3 — S6 — 2025/2026}
\fancyfoot[C]{\small\thepage}

% ─────────────────────────────────────────────────────────────────────────────
\begin{document}
\thispagestyle{empty}

% ══ HEADER PAGE 1 ══════════════════════════════════════════════════════════════
\noindent
\begin{minipage}[c]{0.42\textwidth}
  \textcolor{uirblue}{\Large\bfseries U\={I}R}\\
  \textcolor{gray}{\footnotesize COLLEGE OF ENGINEERING \& ARCHITECTURE}
\end{minipage}%
\hfill
\begin{minipage}[c]{0.52\textwidth}
  \raggedleft
  \textcolor{uirblue}{\textbf{Ecole Supérieure}}\\
  \textcolor{uirblue}{\textbf{d'Informatique et du Numérique}}\\
  \textcolor{gray}{\footnotesize COLLEGE OF ENGINEERING \& ARCHITECTURE}
\end{minipage}

\vspace{6pt}
\noindent\rule{\linewidth}{0.8pt}
\vspace{4pt}

\begin{center}
  \textit{Examen Final (CF)}\\[6pt]
  {\LARGE\bfseries ESIN : Semestre 6}
\end{center}

\vspace{6pt}
\noindent
\begin{tabular}{ll}
  \textbf{Matière}     & Web Development I \\
  \textbf{Enseignant}  & Saad NOUFEL \\
  \textbf{Date}        & Mai 2026 \\
  \textbf{Durée}       & 2h00 \\
\end{tabular}

\vspace{8pt}

% ── Course Outcomes table ─────────────────────────────────────────────────────
\begin{center}
\begin{tabular}{|>{\bfseries}p{10cm}|c|c|}
  \hline
  \rowcolor{uirblue!15}
  \textbf{Course Outcomes} & \textbf{Exercice} & \textbf{Pourcentage} \\
  \hline
  CO2 : Mastery of Basic and OOP principles in PHP to develop robust and scalable applications
    & 3 & 45\% \\
  \hline
  CO3 : Database access with PDO and web authentication via Sessions \& Cookies
    & 1 \& 2 & 55\% \\
  \hline
\end{tabular}
\end{center}

\vspace{10pt}

% ── Student info ──────────────────────────────────────────────────────────────
\noindent
\begin{tabular}{ll}
  Nom              & \dotfill \\[8pt]
  Prénom           & \dotfill \\[8pt]
  Numéro d'étudiant & \dotfill \\
\end{tabular}

\vspace{10pt}

% ── Consignes ─────────────────────────────────────────────────────────────────
\noindent\textbf{Consignes :}
\begin{itemize}[leftmargin=1.5em, itemsep=1pt, topsep=2pt]
  \item Répondez directement sur le sujet
  \item Vous disposez d'un cheatsheet en dernière page
  \item \textbf{Bonne chance !}
\end{itemize}

\vspace{6pt}
\noindent\rule{\linewidth}{0.4pt}

% ══════════════════════════════════════════════════════════════════════════════
%  EXERCICE 1 — Sessions & Cookies
% ══════════════════════════════════════════════════════════════════════════════
\begin{mdframed}[style=exstyle, frametitle={Exercice 1 \quad Sessions \& Cookies \hfill (5 pts)}]
On souhaite développer un système d'authentification pour \textbf{CampusMedia}, une vidéothèque universitaire en ligne.

Soit le formulaire HTML de connexion suivant :

\begin{lstlisting}[language=HTML2]
<form method="POST" action="connexion.php">
    <label>Login :</label>
    <input type="text" name="login" required>

    <label>Mot de passe :</label>
    <input type="password" name="mot_de_passe" required>

    <button type="submit">Se connecter</button>
</form>
\end{lstlisting}
\end{mdframed}

\bigskip

\noindent\textbf{Question 1 (2 pts)} : Écrivez le code PHP du fichier \texttt{connexion.php} qui :
\begin{itemize}[itemsep=1pt, topsep=2pt]
  \item Démarre la session
  \item Récupère le login et le mot de passe envoyés par le formulaire
  \item Vérifie que les deux champs ne sont pas vides, sinon affiche \texttt{"Veuillez remplir tous les champs."}
  \item Interroge la table \texttt{utilisateurs} avec PDO pour trouver l'utilisateur dont le \texttt{login} correspond
  \item Vérifie le mot de passe avec \texttt{password\_verify()}
  \item En cas de succès : stocke le login dans \texttt{\$\_SESSION['utilisateur']} et le rôle dans \texttt{\$\_SESSION['role']}, crée un cookie \texttt{derniere\_connexion} avec la date et l'heure actuelles (valide \textbf{30 jours}), redirige vers \texttt{catalogue.php}
  \item En cas d'échec : affiche \texttt{"Identifiants incorrects."}
\end{itemize}

\answerbox{130pt}

\noindent\textbf{Question 2 (1.5 pts)} : Écrivez le code PHP du \textbf{début} du fichier \texttt{catalogue.php} qui :
\begin{itemize}[itemsep=1pt, topsep=2pt]
  \item Vérifie que l'utilisateur est connecté, sinon redirige vers \texttt{connexion.php}
  \item Affiche \texttt{"Bienvenue, [login] ! (rôle : [role])"} en récupérant les données depuis la session
  \item Si le cookie \texttt{derniere\_connexion} existe, affiche \texttt{"Dernière connexion : [date]"}, sinon \texttt{"Première connexion"}
\end{itemize}

\answerbox{90pt}

\noindent\textbf{Question 3 (1.5 pts)} : Écrivez le code complet du fichier \texttt{deconnexion.php} qui :
\begin{itemize}[itemsep=1pt, topsep=2pt]
  \item Démarre la session, supprime toutes les variables de session et détruit la session côté serveur
  \item Supprime le cookie \texttt{derniere\_connexion}
  \item Redirige vers \texttt{connexion.php}
\end{itemize}

\answerbox{80pt}

\newpage

% ══════════════════════════════════════════════════════════════════════════════
%  EXERCICE 2 — PHP et PDO
% ══════════════════════════════════════════════════════════════════════════════
\begin{mdframed}[style=exstyle, frametitle={Exercice 2 \quad PHP et PDO \hfill (6 pts)}]
On souhaite gérer le catalogue de films de \textbf{CampusMedia}.
Les films sont stockés dans la table \texttt{FILMS} de la base de données \texttt{campusmedia\_bd}.

\vspace{6pt}
\noindent Structure de la table \texttt{FILMS} :

\vspace{4pt}
\noindent
\begin{tabular}{|l|l|}
  \hline
  \rowcolor{uirblue!10}
  \textbf{Colonne} & \textbf{Type} \\
  \hline
  \texttt{id}          & \texttt{INT AUTO\_INCREMENT PRIMARY KEY} \\
  \texttt{titre}       & \texttt{VARCHAR(100)} \\
  \texttt{realisateur} & \texttt{VARCHAR(80)} \\
  \texttt{genre}       & \texttt{VARCHAR(30)} \\
  \texttt{annee}       & \texttt{YEAR} \\
  \texttt{note}        & \texttt{DECIMAL(3,1)} \\
  \hline
\end{tabular}

\vspace{6pt}
\noindent Informations de connexion : \texttt{mysql:host=localhost;dbname=campusmedia\_bd},
utilisateur : \texttt{'mediauser'}, mot de passe : \texttt{'media2026'}
\end{mdframed}

\bigskip

\noindent\textbf{Question 1 (1 pt)} : Écrivez la fonction \texttt{getConnexion(): PDO} permettant de se connecter à la base de données avec le mode d'erreur \texttt{ERRMODE\_EXCEPTION} et le mode de fetch \texttt{FETCH\_ASSOC} par défaut.

\answerbox{80pt}

\noindent\textbf{Question 2 (2 pts)} : Écrivez la méthode statique \texttt{getAll(PDO \$pdo, int \$limit = 0): array} qui :
\begin{itemize}[itemsep=1pt, topsep=2pt]
  \item Récupère tous les films depuis la base de données, triés par note décroissante
  \item Si \texttt{\$limit > 0}, limite le nombre de résultats retournés
  \item Retourne un tableau d'objets \texttt{Film}
\end{itemize}

\answerbox{110pt}

\noindent\textbf{Question 3 (1.5 pts)} : Écrivez la méthode \texttt{save(PDO \$pdo): bool} qui insère le film courant (\texttt{\$this}) dans la base de données en utilisant des \textbf{paramètres nommés} (\texttt{:label}). En cas de succès, mettez à jour \texttt{\$this->id} avec l'identifiant généré.

\answerbox{110pt}

\noindent\textbf{Question 4 (1.5 pts)} : Écrivez la méthode statique \texttt{findByGenre(PDO \$pdo, string \$genre): array} qui retourne tous les films d'un genre donné triés par note décroissante. Expliquez ensuite en une phrase pourquoi votre implémentation protège contre les injections SQL.

\answerbox{90pt}

\noindent\textit{Explication} : \dotfill

\vspace{4pt}
\dotfill

\newpage

% ══════════════════════════════════════════════════════════════════════════════
%  EXERCICE 3 — POO
% ══════════════════════════════════════════════════════════════════════════════
\begin{mdframed}[style=exstyle, frametitle={Exercice 3 \quad POO en PHP \hfill (9 pts)}]
On souhaite modéliser les ressources disponibles dans la vidéothèque \textbf{CampusMedia}.
\end{mdframed}

\bigskip

\noindent\textbf{Question 1 (2 pts)} : Créez une classe \textbf{abstraite} \texttt{Ressource} avec :
\begin{itemize}[itemsep=1pt, topsep=2pt]
  \item Attributs \textbf{protégés} : \texttt{id} (int), \texttt{titre} (string), \texttt{dateAjout} (string)
  \item Un constructeur pour initialiser ces trois attributs
  \item Une méthode \textbf{abstraite} \texttt{getType(): string}
  \item Une méthode \textbf{concrète} \texttt{\_\_toString()} qui retourne une chaîne contenant l'id, le titre et la date d'ajout
\end{itemize}

\answerbox{120pt}

\noindent\textbf{Question 2 (1 pt)} : Créez une \textbf{interface} \texttt{Empruntable} avec les méthodes suivantes :
\begin{itemize}[itemsep=1pt, topsep=2pt]
  \item \texttt{emprunter(string \$nomEtudiant): bool} --- retourne \texttt{true} si l'emprunt est réussi, \texttt{false} sinon
  \item \texttt{retourner(): bool} --- retourne \texttt{true} si le retour est réussi, \texttt{false} sinon
  \item \texttt{getEmprunteur(): ?string} --- retourne le nom de l'emprunteur ou \texttt{null} si la ressource est disponible
\end{itemize}

\answerbox{80pt}

\noindent\textbf{Question 3 (6 pts)} : Créez une classe \texttt{Film} qui \textbf{hérite} de \texttt{Ressource} ET \textbf{implémente} \texttt{Empruntable} avec :
\begin{itemize}[itemsep=1pt, topsep=2pt]
  \item Attributs propres : \texttt{realisateur} (string), \texttt{genre} (string), \texttt{dureeMinutes} (int), \texttt{estDisponible} (bool), \texttt{emprunteur} (?string)
  \item Un constructeur qui appelle \texttt{parent::\_\_construct()} et initialise les attributs propres (\texttt{estDisponible = true} et \texttt{emprunteur = null} par défaut)
  \item Un \textbf{getter} pour \texttt{realisateur}
  \item Un \textbf{setter} pour \texttt{genre}
  \item Implémentation de \texttt{getType()}: retourne la chaîne \texttt{"Film"}
  \item Implémentation des \textbf{trois méthodes} de \texttt{Empruntable} :\\
    --- \texttt{emprunter()} : réussi uniquement si le film est disponible\\
    --- \texttt{retourner()} : réussi uniquement si le film est actuellement emprunté\\
    --- \texttt{getEmprunteur()} : retourne l'emprunteur ou \texttt{null}
  \item Une méthode \texttt{\_\_toString()} qui surcharge celle du parent pour inclure également le réalisateur, le genre, la durée et la disponibilité
\end{itemize}

\answerbox{210pt}

\newpage

% ══════════════════════════════════════════════════════════════════════════════
%  QCM
% ══════════════════════════════════════════════════════════════════════════════
\begin{center}
  {\Large\bfseries QCM --- Cochez la ou les bonnes réponses \hspace{1em} (5 pts)}
\end{center}

\vspace{8pt}

\noindent\textbf{Question 1 (1 pt)} : Laquelle des affirmations suivantes est correcte concernant \texttt{session\_start()} ?

\begin{itemize}[leftmargin=2em, itemsep=3pt, topsep=4pt]
  \item[\cb] Elle peut être appelée après l'envoi de contenu HTML
  \item[\cb] Elle doit être appelée avant tout output, en tout premier
  \item[\cb] Elle crée automatiquement un cookie côté client avec l'ID de session
  \item[\cb] Elle stocke les données de session dans le navigateur
\end{itemize}

\vspace{6pt}

\noindent\textbf{Question 2 (1 pt)} : Comment supprimer définitivement un cookie nommé \texttt{token} en PHP ?

\begin{itemize}[leftmargin=2em, itemsep=3pt, topsep=4pt]
  \item[\cb] \texttt{unset(\$\_COOKIE['token'])}
  \item[\cb] \texttt{setcookie('token', '', time() - 3600, '/')}
  \item[\cb] \texttt{deletecookie('token')}
  \item[\cb] \texttt{session\_destroy()}
\end{itemize}

\vspace{6pt}

\noindent\textbf{Question 3 (1 pt)} : Parmi les instructions PDO suivantes, laquelle protège contre les injections SQL ?

\begin{itemize}[leftmargin=2em, itemsep=3pt, topsep=4pt]
  \item[\cb] \texttt{\$pdo->query("SELECT * FROM films WHERE genre = '\$genre'");}
  \item[\cb] \texttt{\$stmt = \$pdo->prepare("SELECT * FROM films WHERE genre = ?"); \$stmt->execute([\$genre]);}
  \item[\cb] \texttt{\$pdo->exec("SELECT * FROM films WHERE genre = " . \$genre);}
  \item[\cb] \texttt{\$pdo->query("SELECT * FROM films WHERE genre = " . htmlspecialchars(\$genre));}
\end{itemize}

\vspace{6pt}

\noindent\textbf{Question 4 (1 pt)} : Lesquelles des affirmations suivantes sur les classes abstraites et les interfaces sont \textbf{vraies} ?

\begin{itemize}[leftmargin=2em, itemsep=3pt, topsep=4pt]
  \item[\cb] Une classe abstraite peut contenir des méthodes avec implémentation
  \item[\cb] Une interface peut contenir des méthodes avec implémentation (corps)
  \item[\cb] Une classe peut hériter de plusieurs classes abstraites à la fois en PHP
  \item[\cb] Une classe peut implémenter plusieurs interfaces à la fois en PHP
\end{itemize}

\vspace{6pt}

\noindent\textbf{Question 5 (1 pt)} : Que retourne \texttt{\$stmt->fetch(PDO::FETCH\_ASSOC)} quand la requête ne trouve aucun résultat ?

\begin{itemize}[leftmargin=2em, itemsep=3pt, topsep=4pt]
  \item[\cb] \texttt{[]}
  \item[\cb] \texttt{null}
  \item[\cb] \texttt{false}
  \item[\cb] Une exception \texttt{PDOException}
\end{itemize}

\vspace{10pt}
\begin{center}
  \rule{6cm}{0.4pt}\\[4pt]
  {\large --- Fin du sujet ---}\\
  \rule{6cm}{0.4pt}
\end{center}

\newpage

% ══════════════════════════════════════════════════════════════════════════════
%  CHEATSHEET
% ══════════════════════════════════════════════════════════════════════════════
\thispagestyle{fancy}
\begin{center}
  {\LARGE\bfseries Web Development I --- Cheatsheet}\\[2pt]
  {\normalsize Chapitres 3, 4 \& 5 --- ESIN Ing3}
\end{center}

\vspace{6pt}
\noindent\rule{\linewidth}{0.8pt}
\vspace{4pt}

\begin{multicols}{2}

% ── LEFT COLUMN ───────────────────────────────────────────────────────────────
\noindent\colorbox{uirblue!15}{\parbox{\linewidth-2pt}{\bfseries CHAPITRE 3 : POO EN PHP}}

\vspace{4pt}
\textcolor{uirblue}{\textbf{Classe abstraite}}
\begin{lstlisting}[language=PHP2, numbers=none]
abstract class Ressource {
    protected int $id;
    public function __construct(int $id) {
        $this->id = $id;
    }
    abstract public function getType(): string;
    public function __toString(): string {
        return "ID: {$this->id}";
    }
}
\end{lstlisting}

\vspace{4pt}
\textcolor{uirblue}{\textbf{Interface}}
\begin{lstlisting}[language=PHP2, numbers=none]
interface Empruntable {
    public function emprunter(string $nom): bool;
    public function retourner(): bool;
    public function getEmprunteur(): ?string;
}
\end{lstlisting}

\vspace{4pt}
\textcolor{uirblue}{\textbf{Héritage + Interface}}
\begin{lstlisting}[language=PHP2, numbers=none]
class Film extends Ressource
             implements Empruntable {
    public function __construct(
        int $id, string $titre
    ) {
        parent::__construct($id, $titre);
    }
    public function getType(): string {
        return "Film";
    }
    public function emprunter(
        string $nom
    ): bool { ... }
}
\end{lstlisting}

\columnbreak

% ── RIGHT COLUMN ──────────────────────────────────────────────────────────────
\noindent\colorbox{uirblue!15}{\parbox{\linewidth-2pt}{\bfseries CHAPITRE 4 : PHP ET PDO}}

\vspace{4pt}
\textcolor{uirblue}{\textbf{Connexion}}
\begin{lstlisting}[language=PHP2, numbers=none]
$pdo = new PDO($dsn, $user, $pass, [
    PDO::ATTR_ERRMODE =>
        PDO::ERRMODE_EXCEPTION,
    PDO::ATTR_DEFAULT_FETCH_MODE =>
        PDO::FETCH_ASSOC,
]);
\end{lstlisting}

\vspace{4pt}
\textcolor{uirblue}{\textbf{Prepared statements}}
\begin{lstlisting}[language=PHP2, numbers=none]
// Param nommes
$stmt = $pdo->prepare(
    "INSERT INTO t (col) VALUES (:col)"
);
$stmt->bindValue(':col', $v, PDO::PARAM_STR);
$stmt->execute();
$id = (int) $pdo->lastInsertId();

// Param positionnels
$stmt = $pdo->prepare(
    "SELECT * FROM t WHERE id = ?"
);
$stmt->execute([$id]);
$row  = $stmt->fetch();     // une ligne
$rows = $stmt->fetchAll();  // toutes
\end{lstlisting}

\vspace{4pt}
\noindent\colorbox{uirblue!15}{\parbox{\linewidth-2pt}{\bfseries CHAPITRE 5 : SESSIONS \& COOKIES}}

\vspace{4pt}
\textcolor{uirblue}{\textbf{Session}}
\begin{lstlisting}[language=PHP2, numbers=none]
session_start(); // avant tout output !
$_SESSION['utilisateur'] = 'admin';
if (!isset($_SESSION['utilisateur'])) {
    header('Location: login.php'); exit;
}
session_unset();
session_destroy(); // deconnexion
\end{lstlisting}

\vspace{4pt}
\textcolor{uirblue}{\textbf{Cookie}}
\begin{lstlisting}[language=PHP2, numbers=none]
// Creer (30 jours)
setcookie('nom', 'val',
          time() + 30*24*3600, '/');
// Lire
echo $_COOKIE['nom'];
// Supprimer
setcookie('nom', '', time()-3600, '/');
\end{lstlisting}

\vspace{4pt}
\textcolor{uirblue}{\textbf{Password + Redirection}}
\begin{lstlisting}[language=PHP2, numbers=none]
$hash = password_hash($mdp, PASSWORD_DEFAULT);
if (password_verify($mdp, $hash)) { ... }

header('Location: page.php'); exit;
\end{lstlisting}

\end{multicols}

\end{document}
